\chapter{Fragmentation et routage}\label{ch:mx:fragmentation-and-routing}

Un ordre d'achat de 5\,000 actions au meilleur cours vendeur trouve 1\,200 titres sur la plateforme à laquelle il a été envoyé. Les 3\,800 autres à ce prix se trouvent sur quatre autres plateformes, et au moment où
un second ordre les atteint, certains ont déjà disparu. Un titre qui se négocie en de nombreux lieux possède de nombreux carnets d'ordres, et aucun d'eux n'est à lui seul le marché~:
le marché est une vue assemblée à partir de l'ensemble, en retard du temps que mettent les données à voyager, et différente pour chaque trader qui l'assemble. Ce
chapitre construit cette vue dans le simulateur de bourse à partir de trois plateformes, compare le flux consolidé aux flux propres des plateformes, compte les \omterm{def:mx:fragmentation-and-routing:locked}{marchés bloqués} et
\omterm{def:mx:fragmentation-and-routing:locked}{croisés}, et se demande dans les cotations de quelle plateforme se forme le prix.

\section{Nombreuses plateformes, un seul titre}

\begin{definition}[Fragmentation des marchés]\label{def:mx:fragmentation-and-routing:fragmentation}
La \emph{fragmentation des marchés}\index{fragmentation des marchés} désigne la négociation d'un même titre sur plusieurs plateformes à la fois, chacune avec son propre carnet d'ordres,
ses propres règles et ses propres participants, de sorte que la liquidité à un prix donné se répartit entre elles.
\end{definition}

Aux États-Unis, une action peut se négocier sur toutes les bourses, sur des systèmes de négociation alternatifs et chez des wholesalers~; en Europe, sur son marché primaire,
dans des systèmes multilatéraux de négociation et chez des internalisateurs systématiques (livre~1, chapitres~4, 9 et 11). La fragmentation permet aux plateformes de rivaliser sur les frais, la rapidité et les types d'ordres,
et offre aux traders des choix~; elle divise aussi la liquidité et fait de chaque ordre une décision de routage. Les données sur son effet net sont mitigées~:
O'Hara et Ye (2011) ont constaté que les actions américaines dont la négociation était plus fragmentée affichaient des coûts de transaction plus faibles, des exécutions plus rapides et des prix plus proches d'une marche aléatoire,
au prix d'une volatilité de court terme plus élevée~; Degryse, de Jong et van Kervel (2015) ont distingué la fragmentation visible de la négociation dans l'ombre sur les actions européennes.

Le marché du chapitre comporte trois plateformes, A, B et C, chacune avec un flux de fond issu de \texttt{firm.tape}. Les trois flux partagent un même prix efficient $P^\ast$,
mais leurs traders le voient à des instants différents~: immédiatement sur A, avec deux secondes de retard sur B et dix secondes sur C, et B et C ont moins d'activité (70~\% et 45~\%
des taux de A). Les plateformes plus tardives représentent des lieux où les cotations sont mises à jour par des participants plus lents. Un arbitragiste qui reçoit des flux directs
à 20 microsecondes des trois plateformes prend les deux côtés chaque fois que le carnet consolidé est \omterm{def:mx:fragmentation-and-routing:locked}{croisé}, et un flux consolidé (\cref{fig:mx:fragmentation-and-routing:views}) publie
les meilleurs cours acheteur et vendeur nationaux (NBBO) une demi-milliseconde après la mise à jour de chaque plateforme. Sur une heure simulée, la taille affichée au meilleur cours vendeur se répartit
en moyenne à 38~\%, 30~\% et 32~\% entre A, B et C~: un ordre envoyé à une plateforme trouve environ un tiers de la taille à ce prix.

\begin{omfigure}
\begin{tikzpicture}[font=\footnotesize,node distance=1.2cm,
  box/.style={draw,rounded corners,minimum width=1.6cm,minimum height=0.7cm,align=center}]
\node[box,fill=omDef!12] (a) at (0,2) {plateforme A};
\node[box,fill=omDef!12] (b) at (0,0.8) {plateforme B};
\node[box,fill=omDef!12] (c) at (0,-0.4) {plateforme C};
\node[box,fill=omProp!15] (sip) at (4.2,-1.6) {flux consolidé\\ (SIP)};
\node[box,fill=omThm!12] (t) at (8.4,0.8) {trader};
\foreach \v in {a,b,c}{\draw[->,thick,omThm] (\v.east) -- (t.west);}
\foreach \v in {a,b,c}{\draw[->,gray] (\v.south east) -- (sip.west);}
\draw[->,gray] (sip.east) -- (t.south west);
\node[omThm,above] at (4.2,1.9) {flux directs~: 20 $\mu$s};
\node[gray,below] at (1.9,-1.3) {0.5 ms};
\node[gray,below right] at (6.1,-1.2) {+ 20 $\mu$s};
\end{tikzpicture}

\omcaption{Deux vues d'un même marché. Un trader disposant de flux directs construit les meilleurs cours à partir des données propres à chaque plateforme~; le flux consolidé reçoit
le meilleur niveau de chaque plateforme, calcule le NBBO et le publie plus tard. Les délais sont ceux de la simulation.}\label{fig:mx:fragmentation-and-routing:views}
\end{omfigure}

\section{Vues consolidée et directe du marché}

Un processeur d'information sur les titres recueille les cotations et les transactions de chaque plateforme et publie le NBBO~; un trader qui s'abonne aux flux directs des plateformes construit le même NBBO
plus tôt (livre~1, chapitre~9). \texttt{firm.consolidated} reconstruit les deux vues à partir des meilleurs niveaux des plateformes, avec le
\texttt{firm\_nbbo} du livre~1 (une cotation n'entre que pour un lot standard), et \texttt{firm.exchsim} exécute désormais un flux consolidé en direct que reçoivent les agents NBBO
(\texttt{SipConfig}, \texttt{Agent.on\_nbbo}).

Avec un flux d'une demi-milliseconde, les deux vues diffèrent pendant 0.0098~\% du temps de l'heure simulée~: le NBBO direct change de prix 0.44 fois par seconde, et
chaque changement n'est vu qu'avec une demi-milliseconde de retard. La part croît avec le délai~: 0.002~\% à 50 microsecondes, 0.018~\% à une milliseconde, 0.084~\% à cinq,
0.33~\% à vingt et 1.5~\% à cent (\cref{fig:mx:fragmentation-and-routing:delay}), en deçà du produit du taux de changement par le délai, car les changements arrivent par rafales. Les parts de temps masquent l'essentiel.
Au moment d'une transaction, les vues diffèrent beaucoup plus souvent, parce que les transactions suivent les changements de cotation~:
le NBBO consolidé juste avant une transaction était périmé pour 18.5~\% des transactions de l'heure, presque toutes celles de l'arbitragiste, qui se déclenchent dans les 40
microsecondes suivant le changement qui les provoque~; pour les transactions du fond, la part est de 0.02~\%. Ding, Hanna et Hendershott (2014) ont mesuré l'écart réel~:
dans les actions très actives, les NBBO consolidé et direct se contredisent plusieurs fois par seconde, pendant une à deux millisecondes à chaque fois.

\begin{omfigure}
\begin{tikzpicture}
\begin{axis}[width=9cm,height=5.4cm,xmode=log,ymode=log,xlabel={\small délai du flux consolidé (ms)},
  ylabel={\small temps avec un NBBO différent (\%)},tick label style={font=\footnotesize},grid=major,grid style={gray!25},
  legend columns=2,legend style={font=\footnotesize,at={(0.5,-0.3)},anchor=north,draw=none,fill=none,
  /tikz/every even column/.append style={column sep=8pt}},table/col sep=comma]
\addplot[omDef,very thick,mark=*] table[x=ms,y=share]{figdata/microstructure/08-fragmentation-and-routing/delay.csv};
\addplot[gray,dashed] table[x=ms,y=linear]{figdata/microstructure/08-fragmentation-and-routing/delay.csv};
\legend{simulé,changements de prix par seconde $\times$ délai}
\end{axis}
\end{tikzpicture}

\omcaption{Fréquence à laquelle une vue consolidée diverge de la vue directe, en fonction de son délai, sur l'heure simulée à trois plateformes. La ligne en pointillés suppose
des changements de prix isolés~; les changements qui arrivent par rafales se chevauchent et sont manqués ensemble. Données~: \texttt{mx\_frag.study}.}\label{fig:mx:fragmentation-and-routing:delay}
\end{omfigure}

\section{Cotations protégées, marchés bloqués et croisés}

\begin{definition}[Cotation protégée]\label{def:mx:fragmentation-and-routing:protected}
Une \emph{cotation protégée}\index{cotation protégée} est un meilleur cours acheteur ou vendeur affiché que les autres plateformes ne peuvent pas traverser~: un ordre exécuté
ailleurs à un prix moins favorable doit d'abord la prendre, ou être routé vers elle. Aux États-Unis, il s'agit de la cotation protégée de la Rule 611 (livre~1,
chapitre~9)~; dans le simulateur et dans \texttt{firm\_nbbo}, une cotation d'au moins un lot standard.
\end{definition}

\begin{definition}[Marchés bloqués et croisés]\label{def:mx:fragmentation-and-routing:locked}
Le marché est \emph{bloqué}\index{marché bloqué} lorsque le meilleur cours acheteur d'une plateforme est égal au meilleur cours vendeur d'une autre, et
\emph{croisé}\index{marché croisé} lorsqu'il lui est supérieur. Aucun des deux ne peut survenir à l'intérieur d'un même carnet, dont le \omterm{def:mx:the-limit-order-book:engine}{moteur d'appariement} exécuterait les deux côtés~; entre plateformes, chaque carnet ne voit que lui-même.
\end{definition}

Un \omterm{def:mx:fragmentation-and-routing:locked}{marché croisé} est un gain sans risque pour celui qui prend les deux côtés~; un \omterm{def:mx:fragmentation-and-routing:locked}{marché bloqué} ne l'est pas (il n'y a rien à gagner à prix égaux, avant frais), mais il
dit aux routeurs deux choses différentes sur le prix. Les règles américaines maintiennent les deux situations rares~: les plateformes ne doivent pas afficher de cotations qui bloquent ou croisent les \omterm{def:mx:fragmentation-and-routing:protected}{cotations protégées} ailleurs, et
les routeurs ne doivent pas traverser les \omterm{def:mx:fragmentation-and-routing:protected}{cotations protégées}. Les plateformes du simulateur n'ont aucune règle de ce type, ce qui montre ce que ces règles empêchent.

\omcode{code/microstructure/08-fragmentation-and-routing/python/mx_frag.py}{52}{69}{L'arbitragiste entre plateformes~: le carnet consolidé issu de ses flux directs, et les deux côtés de tout croisement pris simultanément.}\label{lst:mx:fragmentation-and-routing:arb}

Sans l'arbitragiste, le marché à trois plateformes
est \omterm{def:mx:fragmentation-and-routing:locked}{bloqué} 34.4~\% de l'heure et \omterm{def:mx:fragmentation-and-routing:locked}{croisé} 10.0~\%, en 69 épisodes d'une durée médiane de 4.8 secondes~: les cotations de la plateforme tardive se placent du mauvais côté
de celles des autres jusqu'à ce que ses traders apprennent le nouveau prix. L'arbitragiste supprime les croisements, 565 d'entre eux d'une durée de vie médiane de 41 microsecondes, et
la part croisée tombe à 0.07~\%, mais il laisse les blocages, qui montent à 42.6~\% avec une médiane de 0.44 seconde. Il négocie 1\,410 fois, surtout contre les cotations périmées de C~:
le volume de C passe de 128\,100 à 230\,900 actions. Prendre une cotation en retard pour apprendre le prix relève de l'arbitrage de latence (livre~11,
chapitre~9)~; Bartlett et McCrary (2019) ont mesuré, à l'aide d'horodatages à la microseconde, à quelle fréquence les traders ont pâti du délai du flux consolidé.

\begin{dated}{2026-09}{Protection des ordres en révision}\label{dat:mx:fragmentation-and-routing:opr}
Le 11 juin 2026, la SEC a proposé d'abroger la Rule 611 de la Regulation NMS, l'interdiction de traversée, ainsi que la Rule 610(e), les restrictions sur le blocage et
le croisement des cotations, avec une période de commentaires de 60 jours après la publication au Federal Register. Le marché simulé sans l'une ni l'autre règle ne montre qu'un côté
de la question~: les blocages et les croisements sont laissés aux arbitragistes, qui suppriment les croisements en quelques microsecondes et les blocages jamais.
\end{dated}

\section{Latence entre les plateformes}

Les données du flux consolidé arrivent en retard parce qu'elles voyagent~: de chaque plateforme vers le processeur, puis dans le calcul des NBBO, puis de nouveau vers l'extérieur. Un trader
colocalisé avec chaque plateforme et abonné à leurs flux directs voit chaque changement plus tôt, de toute la différence. Cet écart ne concerne pas
seulement les traders rapides. Un routeur qui envoie un ordre vers la plateforme que le flux consolidé affiche au meilleur cours l'envoie à un prix qui peut
avoir disparu~; un ordre immédiatement exécutable, contrôlé par rapport à un NBBO périmé, peut provoquer un trade-through d'un meilleur prix que le routeur ne voyait pas, ou
le manquer. Le routage (chapitre~18) part d'une vision du marché, et cette vision a un âge~: chaque NBBO qu'un routeur utilise doit porter l'heure qu'il reflète.

\section{Où le prix est-il découvert~?}

\begin{definition}[Formation des prix]\label{def:mx:fragmentation-and-routing:discovery}
La \emph{formation des prix}\index{formation des prix} est le processus par lequel une nouvelle information est intégrée dans les prix. Lorsqu'un même titre est négocié sur plusieurs
marchés, leurs prix partagent un prix efficient commun, et la question est de savoir quel marché fait d'abord bouger ce prix par ses cotations.
\end{definition}

Les prix $p_{j,t}$ des plateformes sont cointégrés~: chacun est le prix efficient commun plus une erreur stationnaire, de sorte que les différences $p_{1,t}-p_{j,t}$ sont
stationnaires (livre~4, chapitre~20). Un modèle vectoriel à correction d'erreur capture à la fois la dynamique de court terme et le rappel vers un prix unique~:
\[
\Delta p_t=\alpha\,\beta' p_{t-1}+\sum_{k=1}^{L}\Gamma_k\,\Delta p_{t-k}+e_t,\qquad \operatorname{Var}(e_t)=\Omega,
\]
avec $\beta'p_{t-1}$ les $n-1$ différences de prix. L'effet permanent d'une innovation sur le prix commun est $\psi e_t$, avec $\psi$ la ligne commune de la matrice
d'impact de long terme, $\psi\propto\alpha_\perp'$ normalisé par $\alpha_\perp'(I-\sum_k\Gamma_k)\iota$.

\begin{definition}[Part d'information]\label{def:mx:fragmentation-and-routing:is}
La \emph{part d'information}\index{part d'information} d'un marché (Hasbrouck, 1995) est la contribution proportionnelle de ses innovations à la variance de
l'innovation du prix efficient commun~: avec $F$ le facteur de Cholesky inférieur de $\Omega$, $\mathrm{IS}_j=([\psi F]_j)^2/(\psi\Omega\psi')$. Le
facteur de Cholesky dépend de l'ordre des marchés lorsque leurs innovations sont corrélées~; les parts sont donc données sous forme de bornes sur l'ensemble des ordres.
\end{definition}

\begin{definition}[Part de composante]\label{def:mx:fragmentation-and-routing:cs}
La \emph{part de composante}\index{part de composante} d'un marché (Gonzalo et Granger, 1995) est son poids dans la composante permanente commune,
$\mathrm{CS}_j=\alpha_{\perp,j}/\sum_k\alpha_{\perp,k}$~: elle est grande pour un marché qui ne s'ajuste pas aux autres, c'est-à-dire pour un marché qui mène.
\end{definition}

\omcode{code/firm/consolidated/firm_consolidated.py}{103}{120}{Parts d'information, bornées sur l'ensemble des ordres, et parts de composante issues du modèle à correction d'erreur.}\label{lst:mx:fragmentation-and-routing:is}

Sur les trois plateformes simulées, avec des cotations au milieu de la fourchette échantillonnées chaque seconde et cinq retards, la \omterm{def:mx:fragmentation-and-routing:is}{part d'information} de la plateforme A se situe entre 0.63 et 0.79, celle de B entre
0.19 et 0.34, celle de C entre 0.01 et 0.04 (\cref{fig:mx:fragmentation-and-routing:shares})~; avec vingt retards, elle va de 0.53 à 0.69, de 0.30 à 0.45 et de 0 à 0.02. Les
\omterm{def:mx:fragmentation-and-routing:cs}{parts de composante} valent 0.75, 0.44 et $-0.20$~: la seconde mesure s'accorde sur le meneur et montre sa faiblesse connue, une part négative pour un marché qui
dépasse l'ajustement. L'application de Hasbrouck lui-même a donné une part médiane de 92.7~\% pour la New York Stock Exchange sur les trente valeurs du Dow.

Les mesures ont besoin que l'avance du meneur soit visible à l'intervalle d'échantillonnage. Avec des délais de 0.2 et 1 seconde au lieu de 2 et 10, les mêmes
estimations placent B en tête (0.37 à 0.59) et A en second (0.23 à 0.44)~: les cotations des plateformes mettent plus de temps que les délais à suivre $P^\ast$, et leur
propre bruit décide du classement. Une part du volume n'est pas non plus une part de la \omterm{def:mx:fragmentation-and-routing:discovery}{formation des prix}~: le volume de C a presque doublé avec l'arbitragiste alors que
sa \omterm{def:mx:fragmentation-and-routing:is}{part d'information} est restée proche de zéro.

\begin{omfigure}
\begin{tikzpicture}
\begin{axis}[width=9.5cm,height=5.2cm,xtick={0,1,2},xticklabels={A,B,C},xmin=-0.5,xmax=2.5,ymin=-0.3,ymax=1,
  ylabel={\small part},tick label style={font=\footnotesize},grid=major,grid style={gray!25},legend columns=1,legend cell align=left,
  legend style={font=\footnotesize,at={(0.5,-0.2)},anchor=north,draw=none,fill=none,/tikz/every even column/.append style={column sep=8pt}},
  table/col sep=comma]
\addplot[omDef,only marks,mark=*,error bars/.cd,y dir=both,y explicit] table[x=k,y=mid,y error=err]{figdata/microstructure/08-fragmentation-and-routing/shares.csv};
\addplot[omProp,only marks,mark=diamond*,mark size=3pt] table[x=k,y=cs]{figdata/microstructure/08-fragmentation-and-routing/shares.csv};
\addplot[omThm,only marks,mark=square*,error bars/.cd,y dir=both,y explicit] table[x expr=\thisrow{k}+0.2,y=mid_close,y error=err_close]{figdata/microstructure/08-fragmentation-and-routing/shares.csv};
\draw[gray] (axis cs:-0.5,0) -- (axis cs:2.5,0);
\legend{{part d'information, délais 0, 2 et 10 s},{part de composante, idem},{part d'information, délais 0, 0.2 et 1 s}}
\end{axis}
\end{tikzpicture}

\omcaption{\omterm{def:mx:fragmentation-and-routing:discovery}{Formation des prix} sur trois plateformes simulées~: bornes des \omterm{def:mx:fragmentation-and-routing:is}{parts d'information} (barres) et \omterm{def:mx:fragmentation-and-routing:cs}{parts de composante}, à partir de cotations au milieu de la fourchette relevées chaque seconde et de cinq retards.
Avec des délais assez grands pour être visibles à l'intervalle d'échantillonnage, les deux mesures trouvent la plateforme qui apprend en premier~; avec des délais courts, elles ne le font pas. Données~:
\texttt{mx\_frag.study}.}\label{fig:mx:fragmentation-and-routing:shares}
\end{omfigure}

\section{Tutoriel~: où le prix est-il trouvé~?}\label{tut:mx:fragmentation-and-routing:tut}

\begin{tutorial}
\textbf{Objectif.} Construire un marché à trois plateformes avec un flux consolidé, comparer les deux visions et mesurer la \omterm{def:mx:fragmentation-and-routing:discovery}{formation des prix}.
\textbf{État final~:} les \cref{fig:mx:fragmentation-and-routing:delay,fig:mx:fragmentation-and-routing:shares} et les chiffres des sections 1 à 5.
\begin{tutsteps}
\item \textbf{Marché.} \texttt{mx\_frag.market(seconds, seed, arb, sip\_ns, lags)}~: trois \texttt{TapeBackground}s sur un même chemin efficient
(\texttt{v\_path}), l'\texttt{Arbitrageur} et \texttt{Simulator(\dots, sip=SipConfig(default\_ns=sip\_ns))}.
\item \textbf{Vues.} \texttt{views(res)}~: les meilleurs cours de chaque plateforme, le NBBO direct (\texttt{firm\_consolidated.nbbo\_events}) et le flux publié
(\texttt{res.sip(1)}).
\item \textbf{Désaccord.} \texttt{disagreement(direct, delayed, lo, hi)} pour des délais de 50 microsecondes à 100 millisecondes~;
\texttt{locked\_crossed} avec et sans l'arbitragiste.
\item \textbf{Formation.} \texttt{information\_shares(mids(tops, grid), lags)} pour deux jeux de délais~; tracer avec \texttt{fig\_frag.py}.
\end{tutsteps}
\textbf{Pour aller plus loin.} Donner à l'arbitragiste un flux à 1 milliseconde et observer la part croisée et le volume de C~; ajouter une règle qui réajuste le prix des ordres
qui bloqueraient ou croiseraient la \omterm{def:mx:fragmentation-and-routing:protected}{cotation protégée} d'une autre plateforme, et mesurer à nouveau les blocages.
\end{tutorial}

\section{Construction~: le marché consolidé}\label{bld:mx:fragmentation-and-routing:consolidated}

\begin{build}
\bfield{Rôle} Un marché unique issu de nombreuses plateformes~: la vision de départ de chaque routeur (chapitre~18), de l'indice de référence des coûts de transaction (chapitre~19) et de la
stratégie inter-plateformes du livre~11, ainsi que des mesures de \omterm{def:mx:fragmentation-and-routing:discovery}{formation des prix}.
\bfield{Interface} \texttt{nbbo\_events(tops, delays, round\_lot)}, \texttt{sample}, \texttt{disagreement}, \texttt{locked\_crossed}~;
\texttt{vecm(p, lags)}, \texttt{information\_shares(p, lags)}. Dans \texttt{firm.exchsim}~: \texttt{SipConfig}, \texttt{Result.sip},
\texttt{Agent.on\_nbbo}, \texttt{ctx.nbbo}, \texttt{ctx.direct\_nbbo}, \texttt{TapeBackground(v\_path=\dots)}.
\bfield{Règles} Le NBBO par \texttt{firm\_nbbo} du livre~1 (lots standard uniquement)~; une vision retardée décale les mises à jour de chaque plateforme de son délai~; vecteurs de cointégration
$p_1-p_j$ imposés~; bornes des \omterm{def:mx:fragmentation-and-routing:is}{parts d'information} sur tous les ordres.
\bfield{Tests d'acceptation} \texttt{code/firm/consolidated/tests/}~: un NBBO construit à la main avec un blocage, une vision retardée et un lot irrégulier~; des \omterm{def:mx:fragmentation-and-routing:is}{parts d'information} qui
trouvent le meneur parmi trois prix simulés décalés de zéro, un et trois pas. \texttt{code/firm/exchsim/tests/}~: les dates de publication et de livraison du flux consolidé
et la vision directe.
\bfield{Pour aller plus loin} Estimation de Johansen des vecteurs de cointégration~; les \omterm{def:mx:fragmentation-and-routing:is}{parts d'information} modifiées qui suppriment le problème de l'ordre~; délais SIP propres
à chaque plateforme, selon la géographie.
\end{build}

\begin{omsources}
\item J.~Hasbrouck, ``One security, many markets: determining the contributions to price discovery'', \emph{Journal of Finance} 50(4), 1995.
\item J.~Gonzalo et C.~Granger, ``Estimation of common long-memory components in cointegrated systems'', \emph{Journal of Business and Economic
Statistics} 13(1), 1995.
\item M.~O'Hara et M.~Ye, ``Is market fragmentation harming market quality?'', \emph{Journal of Financial Economics} 100(3), 2011.
\item S.~Ding, J.~Hanna et T.~Hendershott, ``How slow is the NBBO? A comparison with direct exchange feeds'', \emph{Financial Review} 49(2), 2014.
\item H.~Degryse, F.~de~Jong et V.~van~Kervel, ``The impact of dark trading and visible fragmentation on market quality'', \emph{Review of Finance}
19(4), 2015.
\item R.~P.~Bartlett et J.~McCrary, ``How rigged are stock markets? Evidence from microsecond timestamps'', \emph{Journal of Financial Markets} 45, 2019.
\item U.S. Securities and Exchange Commission, \emph{The Trade-Through Rule and Locked and Crossed Markets}, proposition, Release 34-105655, 2026.
\end{omsources}

\section{Exercices}

\begin{exercise}[$\star$]\label{exo:mx:fragmentation-and-routing:1}
Quatre plateformes cotent~: V1 à 10.00 pour 300 et à 10.02 pour 200~; V2 à 10.01 pour 500 et à 10.03 pour 100~; V3 à 10.01 pour 50 et à 10.02 pour 400~; V4 à 9.99 pour 1\,000 et
10.02 pour 1\,000. Avec des lots standard de 100, quel est le NBBO, et avec quelles tailles~?
\end{exercise}

\begin{exercise}[$\star$]\label{exo:mx:fragmentation-and-routing:2}
La plateforme A affiche un cours acheteur à 10.02 et la plateforme B un cours vendeur à 10.02~; plus tard, A affiche un cours acheteur à 10.03. Nommer les deux états du marché, et dire ce qu'un arbitragiste peut faire dans chacun.
\end{exercise}

\begin{exercise}[$\star$]\label{exo:mx:fragmentation-and-routing:3}
Le cours NBBO d'une action change 500 fois par seconde et le flux consolidé accuse un retard d'une milliseconde. Au plus, quelle part du temps les deux vues peuvent-elles
différer~? Quelle est la même borne pour l'heure simulée avec un retard d'une demi-milliseconde, et pourquoi la part mesurée est-elle inférieure à celle-ci~?
\end{exercise}

\begin{exercise}[$\star\star$]\label{exo:mx:fragmentation-and-routing:4}
Deux marchés ont $\psi=(0.6,0.4)$ et $\Omega=\begin{pmatrix}1&0.5\\0.5&1\end{pmatrix}$. Calculer les bornes de \omterm{def:mx:fragmentation-and-routing:is}{part d'information} de chacun.
\end{exercise}

\begin{exercise}[$\star\star$]\label{exo:mx:fragmentation-and-routing:5}
Dans un modèle à correction d'erreur à deux marchés, avec $z_{t-1}=p_{1,t-1}-p_{2,t-1}$, les coefficients d'ajustement valent $\alpha=(-0.2,0.3)'$. Trouver $\alpha_\perp$ et les
\omterm{def:mx:fragmentation-and-routing:cs}{parts de composante}. Quel marché mène~?
\end{exercise}

\begin{exercise}[$\star\star$]\label{exo:mx:fragmentation-and-routing:6}
La \omterm{def:mx:fragmentation-and-routing:cs}{part de composante} de la plateforme C vaut $-0.20$. Comment une part peut-elle être négative, et que dit-elle des prix de C~?
\end{exercise}

\begin{exercise}[$\star\star\star$]\label{exo:mx:fragmentation-and-routing:7}
\emph{Programmation.} Relancer \texttt{market} avec des retards de 0, 0.2 et 1 seconde, puis estimer les \omterm{def:mx:fragmentation-and-routing:is}{parts d'information} à partir des cotations médianes à une seconde avec cinq
retards. Quelle plateforme vient en premier, et pourquoi la réponse n'est-elle pas celle de la conception~?
\end{exercise}

\begin{exercise}[$\star\star\star$]\label{exo:mx:fragmentation-and-routing:8}
\emph{Trouver la faille.} «~Le volume de la plateforme C a presque doublé ce mois-ci, donc elle contribue désormais beaucoup plus à la \omterm{def:mx:fragmentation-and-routing:discovery}{formation des prix}.~»
\end{exercise}

\section{Problème~: où le prix se forme-t-il~?}

\begin{problem}[{Problème du week-end --- où le prix se forme-t-il~?}]\label{pb:mx:fragmentation-and-routing:1}
Un régulateur demande si un marché à trois plateformes a besoin d'un flux consolidé, d'une règle de protection des ordres et d'une interdiction des blocages. Construire le marché et répondre
par des chiffres.

\medskip\noindent\textbf{Partie I --- Le marché.}
\begin{enumerate}
\item Définir la \omterm{def:mx:fragmentation-and-routing:fragmentation}{fragmentation des marchés}, et énoncer un avantage et un coût.
\item Décrire les plateformes simulées et la manière dont leurs traders apprennent $P^\ast$.
\item Comment la taille affichée au meilleur cours vendeur se répartit-elle entre les plateformes~?
\item Qu'ont trouvé O'Hara et Ye au sujet des actions américaines fragmentées~?
\end{enumerate}

\medskip\noindent\textbf{Partie II --- Deux vues.}
\begin{enumerate}[resume]
\item Comment le flux consolidé construit-il et publie-t-il le NBBO~?
\item Quelle part de l'heure les vues directe et consolidée diffèrent-elles avec un flux d'une demi-milliseconde, et comment cette part croît-elle avec le retard~?
\item Quelle part des transactions a rencontré un NBBO consolidé périmé, et à qui appartenaient-elles~?
\item Qu'ont mesuré Ding, Hanna et Hendershott sur les marchés réels~?
\end{enumerate}

\medskip\noindent\textbf{Partie III --- Blocages et croisements.}
\begin{enumerate}[resume]
\item Définir les \omterm{def:mx:fragmentation-and-routing:locked}{marchés bloqués} et \omterm{def:mx:fragmentation-and-routing:locked}{croisés}.
\item Donner les parts de \omterm{def:mx:fragmentation-and-routing:locked}{marché bloqué} et \omterm{def:mx:fragmentation-and-routing:locked}{croisé} sans l'arbitragiste.
\item Que change l'arbitragiste, et que laisse-t-il inchangé~?
\item Contre les cotations de qui échange-t-il, et comment appelle-t-on cette transaction~?
\item Qu'a proposé la SEC en juin 2026~?
\end{enumerate}

\medskip\noindent\textbf{Partie IV --- Découverte des prix et verdict.}
\begin{enumerate}[resume]
\item Écrire le modèle à correction d'erreur et la \omterm{def:mx:fragmentation-and-routing:is}{part d'information}.
\item Donner les bornes de \omterm{def:mx:fragmentation-and-routing:is}{part d'information} et les \omterm{def:mx:fragmentation-and-routing:cs}{parts de composante} des trois plateformes.
\item Pourquoi les bornes diffèrent-elles d'un ordre de classement à l'autre~?
\item Que se passe-t-il avec des retards de 0.2 et de 1 seconde, et pourquoi~?
\item Énoncer le \emph{résultat nommé}~: les \omterm{def:mx:fragmentation-and-routing:is}{parts d'information} des trois plateformes simulées, et la part du temps pendant laquelle les meilleurs cours du flux direct et du flux consolidé
diffèrent, en fonction du retard du flux consolidé.
\item Le partage du volume mesure-t-il la \omterm{def:mx:fragmentation-and-routing:discovery}{formation des prix}~? Prendre la plateforme C pour exemple.
\item En une phrase~: où le prix se forme-t-il~?
\end{enumerate}
\end{problem}

\section{Questions d'entretien}

\begin{interviewq}[$\star$ \iqroles{trader}]\label{iq:mx:fragmentation-and-routing:1}
Qu'est-ce que le NBBO, qui le calcule, et pourquoi votre propre vision de celui-ci pourrait-elle différer~?
\end{interviewq}

\begin{interviewq}[$\star\star$ \iqroles{trader}]\label{iq:mx:fragmentation-and-routing:2}
Qu'est-ce qu'un \omterm{def:mx:fragmentation-and-routing:locked}{marché bloqué}, et pourquoi une plateforme refuserait-elle d'afficher un ordre qui le bloque~?
\end{interviewq}

\begin{interviewq}[$\star\star$ \iqroles{chercheur}]\label{iq:mx:fragmentation-and-routing:3}
Comment mesurez-vous laquelle de deux plateformes mène la \omterm{def:mx:fragmentation-and-routing:discovery}{formation des prix}~? Qu'est-ce qui peut mal tourner~?
\end{interviewq}

\begin{interviewq}[$\star\star$ \iqroles{développeur}]\label{iq:mx:fragmentation-and-routing:4}
Vous construisez un carnet consolidé à partir de cinq flux directs. Que faut-il à chaque entrée, en plus du prix et de la taille, et pourquoi~?
\end{interviewq}

\begin{interviewq}[$\star\star$ \iqroles{chercheur}]\label{iq:mx:fragmentation-and-routing:5}
La fragmentation est-elle bonne ou mauvaise pour la qualité de marché~? Quelles preuves chercheriez-vous~?
\end{interviewq}

\begin{interviewq}[$\star\star\star$ \iqroles{trader, chercheur}]\label{iq:mx:fragmentation-and-routing:6}
Le flux consolidé accuse un retard d'une demi-milliseconde sur les flux directs. Qui en pâtit, dans quelle mesure, et comment le mesureriez-vous~?
\end{interviewq}
